\sectionTitle{教育背景}{\faiconsixbf{graduation-cap}}

\datedsubsection{\texthl{海南大学} / \textit{计算机科学与技术} / \textit{本科}}{2021.09 -- 2025.06}
\normalline{英语 CET-6 、2022年中国高校计算机大赛——大学生 RPA+AI 创新挑战赛省级二等奖}

\sectionTitle{工作经历}{\faiconsixbf{building-user}} % building-columns

\datedsubsection{\textbf{造物维度美术设计工作室}\ \ 广州}{2024.11 -- 2024.12}
\role{实习}{UI设计助理}
\begin{itemize}
  \item 参与建立可复用组件库，统一图标系统与视觉语言，提升设计一致性。
  \item 参与用户测试与可用性走查，提出7项体验优化点，其中5项被快速落地。
  \item 配合开发输出切图与标注，高效完成5个大版本的设计交付。
\end{itemize}
\datedsubsection{\textbf{蓝船网络科技有限公司}\ \ 广州}{2024.07 -- 2024.08}
\role{远程实习}{}
移动游戏发行商国籍识别与榜单数据采集系统
\begin{itemize}
  \item 独立开发基于 Python + Pandas 的 App Store / Google Play 元数据采集与处理系统。
  \item 设计域名后缀解析 + 地址关键词匹配的国籍推断规则，累计处理 5,000+ 款应用数据，导出标准化 Excel 分析报表。
\end{itemize}

\sectionTitle{项目经历}{\faiconsixbf{users}}

\datedsubsection{\textbf{基于区块链的供应链信用与风险系统}}{2025.01 -- 2025.03}
\role{Blockchain, Python}{个人项目}
设计并实现基于区块链的供应链信用评估系统，实现信用数据上链与不可篡改存储。
\begin{itemize}
  \item 使用 Solidity 开发智能合约，实现信用评分更新与激励机制。
  \item 基于 Flask + Web3 构建后端服务，实现链上链下数据交互。
  \item 使用 React 开发前端信用排行榜系统，实现企业信用可视化。
  \item 基于 PyTorch 构建欺诈检测模型，引入 交易风险控制机制，使异常交易识别能力提升。
  \item 构建完整仿真数据流程，实现信用评分 + 风控闭环。
\end{itemize}

\datedsubsection{\textbf{企业信用风险检测系统}}{2023.03 -- 2023.05}
\role{Java, MySQL}{系统开发}
实现企业工商、财务、司法等多维度信用信息的统一管理，构建自动化风险评分与预警模型，并向业务人员输出可执行的风险处置建议。
\begin{itemize}
  \item 使用 HTML + CSS + JavaScript 完成响应式管理界面开发，包含信用信息看板、风险列表、预警弹窗及数据可视化图表交互，确保用户在不同终端均能流畅操作。
  \item 基于 Java 语言编写业务逻辑层，设计并实现了企业信息管理、定时风险扫描、风险指标计算、预警阈值触发及建议生成等核心模块；通过 Tomcat 容器部署应用，处理前后端 HTTP 请求与数据交互。
  \item 采用 MySQL 数据库设计企业基础信息表、信用评分记录表、风险事件日志表。
\end{itemize}

% Reference Test
% \datedsubsection{\textbf{Paper Title\cite{zaharia2012resilient}}}{May. 2015}
% An xxx optimized for xxx\cite{verma2015large}
% \begin{itemize}
%  \item main contribution
% \end{itemize}

\newpage

\sectionTitle{技能清单}{\faiconsixbf{gears}}

\begin{onehalfspacing}
  \normalline{编程语言与开发框架：熟练使用 C/C++ 及 C++ STL 进行工程开发；熟练使用 Python（Pandas、Flask、Web3）开展数据处理与后端服务开发；具备 Java Web 应用开发与 Tomcat 部署经验；了解 Solidity 智能合约开发。}
  \normalline{数据结构与算法：扎实掌握数组、链表、栈、队列、树、图等数据结构，以及排序、查找、分治、动态规划、贪心等算法；具备复杂度分析能力，能够结合业务场景进行性能优化。}
  \normalline{数据库：熟练使用 SQL 与 MySQL，具备关系型数据库表结构设计、索引优化、查询调优与基础性能排查的实践经验。}
  \normalline{操作系统与 Linux：熟悉进程/线程调度、内存管理、文件系统等核心机制，具备 Linux 环境下的开发、调试与问题定位能力。}
  \normalline{机器学习：具备模型训练与评估实践经验，熟练使用 PyTorch 构建分类、异常检测模型，掌握特征工程与模型评估方法。}
  \normalline{计算机网络与软件工程：熟悉 TCP/IP、HTTP 等网络基础及需求分析、设计模式、测试等工程流程，具备从需求拆解到开发交付的问题定位与项目实施能力。}
  \normalline{其他：持有 C1 驾驶执照。}
\end{onehalfspacing}

\sectionTitle{校园经历}{\faiconsixbf{school}}

\datedsubsection{\textbf{海南大学笛箫协会\ \ 会长}}{2022.06 -- 2023.06}
\begin{onehalfspacing}
\normalline{在职期间，积极负责地开展社团工作，处理社团大小事宜。主持举办了笛箫协会新生见面会、笛箫协会新生才艺挑战赛、笛箫协会第十二届民乐晚会、笛箫知识专家讲座以及其余多项线上线下小型活动。积累了丰富的社团运营与活动策划经验，有效锻炼了团队统筹与协作能力。}
\end{onehalfspacing}

\sectionTitle{自我评价}{\faiconsixbf{comment}} % note-sticky

\begin{onehalfspacing}
\normalline{具备扎实的计算机基础与良好的代码习惯。熟悉 C++ 与 Python 双技术栈，能够利用 STL 高效开发，也能使用 Python 快速完成数据处理与后端服务搭建。在多个项目中独立完成全流程工作，具备一定的全栈思维。乐于学习新技术，善于协作与沟通。团队意识强，能在多元环境中做好跨部门沟通与协同配合。具备出色的学习力与抗压能力，面对新任务能快速上手、稳步推进。为人务实靠谱，将责任心视为职业底线，愿在新平台深耕细作，成为团队中值得信赖的稳定器与贡献者。}
\end{onehalfspacing}
